\section{Relative prime-circle enumeration at the actual prime density}
\label{sec:relativeprimecircles}
The logarithmic-density theorem has an independent arithmetic
consumer on the unweighted prime-sum graph itself.
The row repair must preserve the coefficient of \(\log N\)
in the largest entry; a generic \(\log N\) inverse bound is
too coarse for that purpose.

\begin{lemma}[A diffuseness-sensitive scaling repair]
\label{lem:diffusionscaling}
Let \(Q\) be a symmetric nonnegative bipartite matrix on
\(n\) vertices with equally sized parts, row degrees
\(|d_i-1|\leq\varepsilon\), and \(\max Q_{ij}\leq A/n\),
\(A\geq2\).
Suppose its degree-normalized matrix has nontrivial gap
at least a fixed \(\delta\in(0,1)\).
If \(\varepsilon\log(2A)^2\) is sufficiently small, it has
a stochastic repair
\[
 \widehat P=\diag(1+z)Q\diag(1+z),\qquad
 \|z\|_\infty\leq C_\delta\varepsilon\log(2A).
\]
The repaired kernel differs from \(Q\) in operator norm
by \(O_\delta(\varepsilon\log(2A))\).
\end{lemma}
\begin{proof}
We refine the inverse estimate in Lemma~\ref{lem:scaling}.
For the reversible walk \(R=\diag(d)^{-1}Q\), remove its
constant and parity projections.
The absolute row sum of its \(j\)-th centered power is at
most a fixed constant.
After the first step the squared row norm is at most
\(C A/n\); the fixed spectral gap and Cauchy--Schwarz
therefore give the alternative bound
\(C_\delta\sqrt A(1-\delta)^{j-1}\).
The stationary degrees are between fixed constants, so
these estimates hold also in the stationary weighted norm.
Summing their minimum gives
\[
 \sum_{j\geq0}\min\{C,C_\delta\sqrt A
                                      (1-\delta)^{j-1}\}
                         \leq C_\delta\log(2A).
\]
The constant mode in \((I+R)^{-1}\) costs only \(1/2\);
the parity mode is removed by the same balanced gauge
as in the earlier scaling proof.
Multiplying by \(\diag(d)^{-1}\) proves that the inverse
of \(\diag(d)+Q\), on its compatible subspace, has
infinity norm at most \(C_\delta\log(2A)\).
Changing to the unweighted balanced gauge costs only a
fixed multiple of this bound.

Now solve the exact equation
\[
 (\diag(d)+Q)z=\one-d-z\mathbin{\odot}(Qz)
\]
in the balanced gauge.
The right side is compatible on both parts, including
its quadratic term.
On the ball
\(\|z\|_\infty\leq 2C_\delta\varepsilon\log(2A)\),
the inverse and quadratic Lipschitz bounds give a
contraction provided \(\varepsilon\log(2A)^2\) is small.
The fixed point is positive after adding one and gives
the exact row sums.
The operator perturbation follows from
\(\|\diag(z)Q+Q\diag(z)+\diag(z)Q\diag(z)\|
 \leq C\|z\|_\infty\).
\end{proof}

\paragraph{Formalized diffuseness estimates and positive inverse reference.}
The maintained owner \texttt{HamiltonianMixingResolvent} proves
an exact finite Neumann expansion, a full-tail smoothing estimate,
and its logarithmic cutoff.
\texttt{HamiltonianDiffuseRepairKernel} applies them to the actual
stochastic cross block. If \(T\) is nonnegative and doubly stochastic,
\(T_{ij}\leq a/m\), \(a\geq1\), and
\(\|T-J_m/m\|\leq\rho\in(0,1)\), its exact
constant/parity-corrected kernel \(K\) satisfies
\[
 (1-\rho)\max_i\sum_j|K_{ij}|
 \leq
 2(1-\rho)\left(\frac{\log(2a)}{-\log\rho}+3\right)+1.
\]
Thus the logarithmic cost is independently bounded, rather than
assumed as a hypothesis on the desired scaling.

For a nonnegative nearly stochastic cross block \(T\), let
\(r_i=\sum_jT_{ij}\), \(t_j=\sum_iT_{ij}\), with
\(|r_i-1|,|t_j-1|\leq\varepsilon\in(0,1/2]\).
Put \(c=1+2\varepsilon\) and
\(\tau=\sum_i(c-r_i)=\sum_j(c-t_j)>0\). The auxiliary matrix
\[
 T^*_{ij}=\frac{T_{ij}}c+
               \frac{(c-r_i)(c-t_j)}{c\tau}
\]
is nonnegative and exactly doubly stochastic.
\texttt{HamiltonianStochasticCompletion} proves these identities and
\[
 \max_i\sum_j|T_{ij}-T^*_{ij}|\leq6\varepsilon,\qquad
 \max_j\sum_i|T_{ij}-T^*_{ij}|\leq6\varepsilon,\qquad
 \|T-T^*\|\leq6\varepsilon .
\]
If \(T_{ij}\leq a/m\), it also proves
\[
 T^*_{ij}\leq\frac{a+9\varepsilon}{(1+2\varepsilon)m},
 \qquad
 \|T^*-J_m/m\|\leq\rho+6\varepsilon
 \quad\hbox{when }\|T-J_m/m\|\leq\rho.
\]
The real and complex mass-to-operator estimate is formalized in
\texttt{HamiltonianSchurBudget}.
All thirty-five new declarations in these four owners pass the
focused foundational-axiom audit.
The completion is only an inverse reference; it is not the final
kernel and supplies no assertion about preserving support.
\paragraph{Formalized full repair on the original support.}
The compatible inverse transfer and the nonlinear row equations
are now proved as well.
\texttt{HamiltonianCompatiblePerturbation} retains the exact parity
projection under the small perturbation.
\texttt{HamiltonianDegreeJacobian} proves its compatibility for
the actual diagonal \(\diag(d)\), before normalization.
\texttt{HamiltonianDiffuseJacobianInverse} consumes the positive
stochastic reference above and its entire \(6\varepsilon\) error.
\texttt{HamiltonianDegreeScalingFixedPoint} retains
\(-z\mathbin{\odot}(Qz)\) in the full quadratic contraction.
\texttt{HamiltonianDegreeScaling} uses the already proved exact
conservation law to remove the final parity residual.

The complete finite output is compiled in
\texttt{HamiltonianDiffuseScaling}. For the nearly stochastic
cross block above, set
\[
 g=2\left(\frac{\log(2(a+9\varepsilon))}{-\log\rho}+3\right)
                       +\frac1{1-\rho},\qquad
 r=4\varepsilon g .
\]
For every \(\rho\in(0,1)\), under the explicit hypotheses
\[
 \|T-J_m/m\|\leq\rho-6\varepsilon,\qquad
                    \varepsilon g^2\leq1/256,
\]
there is one vector \(z\) with \(\|z\|_\infty\leq r\)
and every \(1+z_i>0\), for which
\[
 \widehat T_{ij}=(1+z_{A,i})T_{ij}(1+z_{B,j})
\]
is exactly doubly stochastic and nonnegative. The same witness satisfies
\[
 \widehat T_{ij}=0\ \Longleftrightarrow\ T_{ij}=0,\qquad
 \widehat T_{ij}\leq(1+r)^2a/m,\qquad
 \|\widehat T-T\|\leq(2r+r^2)(1+\varepsilon).
\]
All twenty-eight declarations in these six further owners pass
the focused foundational-axiom audit.
For fixed \(\rho\), the bound is
\(g=O_\rho(\log(2a))\), giving the required
\(\varepsilon\log(2a)^2\) smallness scale.
The native degree-normalized spectral consumer of
Lemma~\ref{lem:diffusionscaling} is now formalized in
\texttt{HamiltonianDiffuseSpectralScaling}, with the advertised
logarithmic cap cost. Here are explicit uniform constants.
Put \(\rho=1-\delta/2\) and
\[
 G_\delta=\frac4{-\log\rho}
           +\frac{6+(1-\rho)^{-1}}{\log4},\qquad
 c_\delta=\min\left\{\frac{\delta}{28}(\log4)^2,\,
                    \frac1{256G_\delta^2}\right\},
 \qquad C_\delta=20G_\delta.
\]
For the physical cap \(A/n\), take \(a=A/2\geq1\);
then \(L=\log(4a)=\log(2A)\).
The already proved scalar inverse estimate has a free scale
parameter: applying it at that scale equal to \(a\) gives
\[
 g(a+9\varepsilon,\rho)\leq G_\delta L .
\]
Its scalar hypotheses require only
\(a\leq(1+\varepsilon)a\), which follows from
\(\varepsilon\geq0\); no dimension bound is substituted.

The single smallness hypothesis
\(\varepsilon L^2\leq c_\delta\) first gives
\(\varepsilon\leq\delta/28<1/2\), and then
\(\varepsilon g^2\leq1/256\).
The actual normalized spectral packet, including the two simple
exceptional modes, yields
\[
 \|T-J/m\|\leq1-\delta+8\varepsilon
             \leq\rho-6\varepsilon.
\]
Thus the full nonlinear repair applies to the original matrix.
Its same witness has \(R=4\varepsilon g\leq1/64\), exact
stochastic rows and columns, positive multipliers, identical
support, and
\[
 \|z\|_\infty\leq C_\delta\varepsilon L,\qquad
 \|\widehat T-T\|\leq C_\delta\varepsilon L,\qquad
 \widehat T_{ij}\leq2T_{ij}.
\]
The exact identity
\(\widehat Q-Q=ZQ+QZ+ZQZ\) retains the full quadratic term.
The zero-defect case uses the identity scaling explicitly.
The theorem
\texttt{hamiltonian\_normalized\_spectral\_physical\_cap\_scaling}
has exactly the physical cap \(A/n\) and the cost \(\log(2A)\);
no additional \(\varepsilon\leq1/2\) hypothesis is assumed.
The actual constant-edge prime transfer, uniform degrees and
half-gap are now formalized in Section~\ref{sec:regularfactors}.
\paragraph{Formalized actual prime-density count.}
\texttt{HamiltonianMatrixGeometry} proves that a finite relabeling
preserves the Euclidean operator norm and that the bipartite lift
has exactly the norm of its cross block.
\texttt{HamiltonianPrimeCrossBlock} identifies the actual block
\[
 T_{m,ij}=\frac{\log(2m)}{2m}
     {\bf1}_{\{2i+2j+3\ {\rm prime}\}},
 \qquad 0\leq i,j<m .
\]
Its rows and columns are exactly the corresponding full prime
degrees. Its centered norm is exactly the full constant/parity
operator error, without a dimension factor.

\texttt{HamiltonianPrimeScalingWindow} proves the uniform
smallness window at this density: for every \(B,t>0\), for all
sufficiently large real \(L\), putting
\(\varepsilon=B/\sqrt L\), \(a=L/2\), and \(\rho=2/3\) gives
\[
 0<\varepsilon\leq1/100,\qquad
 \varepsilon g(a+9\varepsilon,\rho)^2\leq t .
\]
The inverse budget is bounded by \(K\log L\) with the explicit
positive constant
\[
 K=\frac{2(1+\log2)}{-\log(2/3)}+9.
\]
The already proved limit \((\log L)^2/\sqrt L\to0\)
closes the condition on an entire terminal interval.

\texttt{HamiltonianPrimeStochasticScaling} applies the full
nonlinear repair to the original block. For every \(\eta>0\),
for every sufficiently large \(m\), there is one positive scaling
with \(\|z\|_\infty\leq\eta\) whose repaired block
\(\widehat T\) is nonnegative, doubly stochastic, and has exactly
the original prime support. The same witness satisfies
\[
 \widehat T_{ij}\leq\frac{(3/5)\log(2m)}m,\qquad
 \|\widehat T-J_m/m\|\leq3/5 .
\]
Finally \texttt{HamiltonianPrimeRelativeCount} consumes these
actual arithmetic inputs in the complete signed counting theorem:
\[
 \left|\left(\frac{m^m}{m!}\right)^2
       m\,\operatorname{cyc}(\operatorname{bip}(\widehat T))-1\right|
                       \leq m^{-3/5}.
\]
All thirty declarations in these five maintained owners pass the
focused foundational-axiom audit.
\paragraph{Formalized entropy and the actual unoriented count.}
\texttt{HamiltonianEntropy} proves the finite supporting-line
inequality, sums every row and column potential exactly, and
proves optimality over every feasible weighting on the original
support, including weightings with zero entries.
\texttt{HamiltonianPrimeEntropy} applies this result to the actual
prime scaling and identifies its entropy exactly.

\texttt{HamiltonianCycleScaling} proves the common vertex-product
weight for every Hamilton cycle before summing the cycles.
\texttt{HamiltonianCycleOrientation} counts actual single-orbit
permutations and their reversal classes. For at least three
vertices, reversal has no fixed point, so the directed count is
exactly twice the unoriented count. Cyclic rotation is already
removed by the permutation representation.

\texttt{HamiltonianPrimeCircleEntropyCount} proves existence and
uniqueness of the finite maximum entropy value \({\cal H}_{2m}\),
and the factorial-normalized relative formula on an entire
terminal interval:
\[
 \left|\left(\frac{m^m}{m!}\right)^2
       (2m)e^{-2{\cal H}_{2m}}h(G_{2m})-1\right|
                                    \leq m^{-3/5}.
\]
The count is the actual finite cardinality of the reversal
classes. \texttt{HamiltonianPrimeCircleAsymptotic} verifies that
its edge relation, under the proved odd/even bijection, is exactly
the relation of distinct integer vertices with prime sum.

\texttt{HamiltonianStirlingNormalization} sums Robbins' exact
stepwise logarithmic bound and takes its established limit.
For every \(m\geq1\), it proves
\[
 \left|\frac{s_m^2}{\pi}-1\right|\leq\frac1{3m},
 \qquad s_m=\frac{m!}{\sqrt{2m}(m/e)^m}.
\]
The full final consumer therefore proves, for every sufficiently
large \(m\),
\[
 \left|\frac{h(G_{2m})}
              {\pi e^{2{\cal H}_{2m}-2m}}-1\right|
                                    \leq 2m^{-3/5}.
\]
All fifty-four declarations in these seven further maintained
owners pass the focused foundational-axiom audit.
Both formulas in Theorem~\ref{thm:relativeprimecircles} now have
complete maintained proofs. The sharp physical repair, full
prime-edge cuts and integral regular-factor criterion with its
actual adjacency loss are now formalized as well in
Section~\ref{sec:regularfactors}. Other advertised arithmetic
consumers and the paper's complete public integration remain
separate formalization obligations.

Let \(G_N\), for even \(N=2m\), have vertices \(1,\ldots,N\)
and edges at distinct prime sums.
For a proper fractional edge weighting \(x\), meaning
\(x_e\geq0\) and \(\sum_{e\ni v}x_e=1\) for every vertex,
write
\[
 {\cal H}_N=\max_x\left(-\sum_{e\in E(G_N)}x_e\log x_e\right).
\]
The optimization is over the actual finite prime-sum graph.
It incorporates its nonuniform degrees and every allowed
edge, with no averaging of the graph itself.

\begin{theorem}[Full relative prime-circle count from its exact entropy]
\label{thm:relativeprimecircles}
Uniformly through even \(N\to\infty\),
\[
 h(G_N)=
 \frac{(m!)^2}{2m\,m^{2m}}\,e^{2{\cal H}_N}
                  \left(1+O(N^{-3/5})\right)
       =\pi e^{\,2{\cal H}_N-N}
                  \left(1+O(N^{-3/5})\right).
\]
Here \(h\) counts unoriented Hamilton cycles modulo cyclic
rotation, as in the earlier prime-circle statements.
The implied constant is absolute eventually.
The formula is unconditional and uses neither RH nor a
prime-pair asymptotic. It is a relative count, rather than
an error of exponential size \(e^{o(N)}\).
\end{theorem}
\begin{proof}
Put \(Q_0=(\log N/N)A_{G_N}\).
Lemma~\ref{lem:unweightedtransfer} supplies the actual
uniform degree defect
\[
 \varepsilon_N=O(\log\log N/\log N)
\]
and operator distance \(o(1)\) from a stochastic kernel
with nontrivial ratio \(1/2+o(1)\).
Consequently the degree normalization of \(Q_0\) also
has nontrivial ratio \(1/2+o(1)\): diagonal normalization
is \(I+O(\varepsilon_N)\), and its constant/parity
vectors differ by \(O(\varepsilon_N)\) from the
normalized constant/parity vectors.
Its largest entry is at most \(\log N/N\).

Apply Lemma~\ref{lem:diffusionscaling} with \(A=\log N\).
The compatibility quantity
\(\varepsilon_N\log(2A)^2=O((\log\log N)^3/\log N)\)
tends to zero. We obtain
\[
 P=\diag(1+z)Q_0\diag(1+z),\qquad
 \|z\|_\infty=O((\log\log N)^2/\log N)=o(1).
\]
Its cross block \(T\), of order \(m=N/2\), is doubly
stochastic and has
\[
 \|T-J_m/m\|\leq1/2+o(1),\qquad
 m\max T_{ij}\leq(1/2+o(1))\log N.
\]
Thus for sufficiently large \(m\) we may take
\(\rho=3/5\), \(A=(3/5)\log(2m)\).
Here
\[
 A-1\leq(3/5)\log m
\]
because \((3/5)\log2\leq1\), so the sharp amplification
consumer of Theorem~\ref{thm:criticalrelative} applies.

The numerical margin is checked explicitly:
\[
 \kappa_{3/5}=\frac{25}{9}\log(17/8)<\frac94,\qquad
 {\cal A}(T)\leq m^{27/20}.
\]
For the strict logarithmic inequality, the first four
terms of \(e^{81/100}\) already exceed \(17/8\).
The two error terms in the comparison are therefore
\(O(\log(2m)^M(m^{-1}+m^{-13/20}))\),
which is \(O(m^{-3/5})\).
The full critical-density signed consumer is now estimated.

Finally set
\(w_v=(1+z_v)\sqrt{\log N/N}\), so that
\(P=\diag(w)A_{G_N}\diag(w)\).
Its positive stochastic edge weighting is the entropy
maximizer. Indeed for every proper \(x\), the supporting
hyperplane inequality for \(-x\log x\) shows optimality
of \(P\), because
\(\log P_{ij}=\log w_i+\log w_j\) on its support and
the vertex sums are fixed.
Strict concavity gives uniqueness of the edge weights;
the harmless bipartite scaling gauge does not change
\(\sum_v\log w_v\).
Hence
\[
 {\cal H}_N=-\sum_v\log w_v.
\]
Every Hamilton cycle has exactly the same \(P\)-weight
\(\prod_v w_v^2=e^{-2{\cal H}_N}\).
Multiplying the relative kernel count by
\(e^{2{\cal H}_N}\) proves the factorial formula.
Stirling gives
\((m!)^2/(2m\,m^{2m})=\pi e^{-N}(1+O(N^{-1}))\),
proving the final normalization.
All arithmetic inputs were already used in the proved
unweighted transfer; the improvement is the retained
signed compensation and its independently bounded remainder.
\end{proof}

\begin{remark}[What precision has changed]
The earlier elementary entropy evaluation gives
\(2{\cal H}_N=N\log N-N\log\log N+o(N)\).
Substituting only that coarse evaluation recovers only
an exponential-base formula.
The theorem above keeps the exact finite variational
quantity and proves a power-saving relative error around
it. It does not silently promote the coarse explicit
expression to a relative formula.
No effective numerical starting order, full all-orders
Filz theorem, or relative classification of every fixed
congruence filter is claimed here.
\end{remark}
